\chapter*{Conclusion générale et perspectives}
\addcontentsline{toc}{chapter}{Conclusion générale et perspectives}
\markboth{Conclusion générale et perspectives}{}

Au terme de la période de développement de ce Projet de Fin d’Études, JOB4YOU est passé d’un simple besoin identifié en tout début de stage à une plateforme fonctionnelle, capable de gérer le processus de recrutement interne de l’entreprise.

Ce travail couvre désormais l’ensemble du cycle de recrutement, de la publication d’une offre jusqu’à la décision finale d’embauche, en passant par la soumission des candidatures, l’analyse automatique des CV par intelligence artificielle, la planification des entretiens et la collecte des feedbacks des managers. La plateforme intègre également deux fonctionnalités différenciantes : le matching inverse « Job idéal », qui suggère à un candidat les offres les plus compatibles avec son CV, et un assistant conversationnel (chatbot) permettant à chacun de suivre ses candidatures en langage naturel — ce dernier étant strictement cloisonné aux données du candidat connecté, à l’exclusion de toute information interne (score IA, notes RH, identité de l’intervieweur). Un système de notifications automatiques par e-mail et un tableau de pilotage visuel de type Kanban complètent l’ensemble et donnent aux équipes RH et aux managers une visibilité en temps réel sur l’avancement de chaque dossier.

Sur le plan technique, ce projet m’a permis d’approfondir mes compétences en développement back-end avec Spring Boot et Spring Security, en développement front-end avec Angular, ainsi qu’en conception de bases de données relationnelles avec PostgreSQL. J’y ai aussi découvert et mis en œuvre des outils d’automatisation (n8n) et d’intelligence artificielle (API de Cohere) au service d’un processus métier concret, tout en renforçant mes compétences en sécurité applicative (authentification JWT, contrôle d’accès basé sur les rôles) et en gestion de projet, selon une démarche itérative inspirée d’Agile/Scrum. J’ai également vérifié le comportement du seuil de décision du scoring sur un jeu de référence annoté (Gold Set) et prévu, pour les cas limites, une validation humaine du score assortie d’une explication automatique des compétences retenues.

Au-delà de l’aspect purement technique, cette expérience au sein d’un grand groupe international m’a beaucoup apporté sur le plan professionnel : travail en équipe, communication régulière avec mon encadrant entreprise, prise en compte de contraintes réelles de sécurité et de fiabilité, et nécessité de concevoir une solution robuste et évolutive dans un contexte professionnel réel.

Plusieurs pistes d’évolution restent ouvertes pour poursuivre ce travail. Côté fonctionnel, la plateforme pourrait accueillir un module d’analyse prédictive du taux de transformation des candidatures, une intégration plus poussée avec des plateformes d’emploi externes (LinkedIn, Indeed), ou un espace candidat plus riche pour un suivi encore plus détaillé de son parcours. La liste des candidatures côté RH expose désormais cette valeur effective (badge du score retenu, mention « corrigé RH », relance de l’analyse depuis une ligne non scorée) ; son tri, exécuté côté client, reste à reporter côté serveur pour qu’il porte sur l’ensemble des dossiers et non sur la seule page affichée. Côté technique, la suite de tests automatisés — qui compte déjà 249 méthodes de test unitaires et de contrôleur (JUnit 5 / Mockito) — gagnerait à être complétée par des tests d’intégration couvrant les services et la base de données, des tests de bout en bout (par exemple Playwright ou Cypress) rejouant les parcours métier complets, et des tests front-end automatisés sur les composants Angular. Le jeu de référence de vingt profils mobilisé pour calibrer le seuil de décision gagnerait à être enrichi de candidatures réelles annotées par les recruteurs, seul moyen d’obtenir une mesure de performance en conditions d’exploitation. Une chaîne d’intégration et de déploiement continus (CI/CD) ainsi qu’un déploiement sur une infrastructure cloud constitueraient des évolutions naturelles pour renforcer la fiabilité et la scalabilité de la plateforme sur le long terme.

Les principales fonctionnalités prévues pour ce projet ont été développées et vérifiées à l’aide des scénarios décrits dans ce rapport : JOB4YOU répond aux besoins identifiés et couvre l’ensemble du cycle, de la publication d’une offre jusqu’à la décision finale. Les objectifs de performance restent en revanche à mesurer, et plusieurs mesures de sécurité et de gouvernance des données doivent être finalisées avant une mise en préproduction — validation du contenu des fichiers téléversés, politique de rétention et d’effacement des données, encadrement des transferts vers Cohere et OpenRouter. Le système constitue donc un prototype fonctionnel, appelé à être industrialisé pour un déploiement chez Sopra HR Software Tunisie, plutôt qu’une solution dont l’aptitude à la production serait déjà démontrée. Cette expérience m’a permis de concevoir, sécuriser, tester et faire évoluer une solution logicielle complète.
